%!TEX program = xelatex
%==========================================================
%  泛函分析讲义 · 中文数学讲义版式（自用）
%==========================================================
\documentclass[a4paper,11pt,oneside,fontset=none,linespread=1.16]{ctexbook}

% ---------- 页面 ----------
\usepackage[a4paper,left=2.1cm,right=2.1cm,top=2.1cm,bottom=2.0cm,
            headsep=0.65cm,footskip=0.9cm]{geometry}

\usepackage{fancyhdr}
\usepackage{amsmath,amssymb,amsthm,mathtools}
\usepackage{mathrsfs}          % \mathscr（花体 B、N 一类）；缺它会报 Undefined control sequence
% 数学字体：Computer Modern（XeLaTeX 默认，CMMI/CMR/CMSY/CMEX）。
% 2026-11-04 定版：不再加载 mathpazo——PazoMath 属 Palatino 家族，希腊字母与
% HTML 版（MathJax TeX / KaTeX_Math）差异明显；改用 CM 后两边字形同族。
% 正文西文仍由下方 \setmainfont{TeX Gyre Pagella} 指定，不受此处影响。
\usepackage[most]{tcolorbox}
\usepackage{xcolor}
\usepackage{etoolbox}
\usepackage{adjustbox}
\usepackage{bm}
\usepackage{enumitem}

\setlist[itemize]{topsep=3pt,itemsep=1pt,parsep=0pt,partopsep=0pt,leftmargin=1.7em}
\setlist[enumerate]{topsep=3pt,itemsep=1pt,parsep=0pt,partopsep=0pt,leftmargin=1.9em}

% ---------- 讲义信息（换章时改这里） ----------
\newcommand{\DocMainTitle}{泛函分析讲义}
\newcommand{\DocSubTitle}{第五章\quad Banach 代数}
\newcommand{\DocSource}{张恭庆、郭懋正《泛函分析讲义（第二版）》下册}
\newcommand{\DocDate}{二〇二六年十一月}
\newcommand{\ShortTitle}{泛函分析讲义}

% ---------- 字体 ----------
\setmainfont{TeX Gyre Pagella}
\IfFontExistsTF{Noto Serif CJK SC}{
    \setCJKmainfont[
        BoldFont = Noto Serif CJK SC Bold,
        ItalicFont = Noto Serif CJK SC Light
    ]{Noto Serif CJK SC}
}{
    \setCJKmainfont[
        BoldFont = FandolSong-Bold.otf,
        ItalicFont = FandolKai-Regular.otf
    ]{FandolSong-Regular.otf}
}
\IfFontExistsTF{Noto Sans CJK SC}{
    \setCJKsansfont[
        BoldFont = Noto Sans CJK SC Bold
    ]{Noto Sans CJK SC}
}{
    \setCJKsansfont[
        BoldFont = FandolHei-Bold.otf
    ]{FandolHei-Regular.otf}
}
\newCJKfontfamily\kaishu{FandolKai-Regular}

% ---------- 颜色 ----------
\definecolor{MainText}{RGB}{25,25,25}
\definecolor{ThmColor}{RGB}{200,80,110}
\definecolor{ThmBg}{RGB}{253,246,248}
\definecolor{DefColor}{RGB}{200,100,30}
\definecolor{DefBg}{RGB}{254,249,243}
\definecolor{ExColor}{RGB}{70,130,180}
\definecolor{ExBg}{RGB}{244,248,252}
\definecolor{ExpColor}{RGB}{106,90,160}
\definecolor{ExpBg}{RGB}{246,244,251}
% 正文内"后补内容"的标注色（深蓝，区别于例题的钢蓝 70,130,180）
\definecolor{AddColor}{RGB}{30,79,148}
\definecolor{ConclColor}{RGB}{18,138,128}
\definecolor{ConclBg}{RGB}{241,249,248}
\color{MainText}

\newcommand{\R}{\mathbb{R}}
\renewcommand{\ge}{\geqslant}
\renewcommand{\geq}{\geqslant}
\renewcommand{\le}{\leqslant}
\renewcommand{\leq}{\leqslant}

% ---------- 标题层级 ----------
\ctexset{
    chapter = {
        format = \centering,
        titleformat = \fontsize{22pt}{27pt}\selectfont\bfseries\color{ThmColor},
        beforeskip = 2pt,
        afterskip = 24pt,
    },
    section = {
        format = \raggedright\Large\bfseries\color{MainText},
        beforeskip = 16pt,
        afterskip = 9pt,
    },
    subsection = {
        format = \raggedright\large\bfseries\color{MainText},
        beforeskip = 12pt,
        afterskip = 6pt,
    },
    subsubsection = {
        format = \raggedright\normalsize\bfseries\color{MainText},
        beforeskip = 10pt,
        afterskip = 4pt,
    }
}

% ---------- 盒子：大盒（定理族 / 定义 / 例题） ----------
\tcbset{
    bigbox/.style={
        enhanced,
        breakable,
        boxrule=0pt,
        sharp corners,
        left=4.0mm, right=3.0mm, top=1.6mm, bottom=2.0mm,
        toptitle=1.0mm, bottomtitle=1.0mm,
        coltitle=white,
        fonttitle=\bfseries\normalsize,
        fontupper=\linespread{1.26}\selectfont,
        before skip=0.85em, after skip=0.55em,
    },
    % 夹注（解释 / 结论）：左侧一条色条，不起标题带
    noteinline/.style={
        enhanced,
        breakable,
        boxrule=0pt,
        sharp corners,
        left=3.2mm, right=2.6mm, top=1.2mm, bottom=1.3mm,
        fontupper=\small\linespread{1.20}\selectfont,
        parskip=0.35em,          % 夹注内多段之间留气口，否则几段挤成一块
        before skip=0.48em, after skip=0.32em,
    }
}

\newtcolorbox{thmBox}[3][]{
    bigbox, colback=ThmBg, colframe=ThmColor,
    left=4.6mm, top=1.9mm, bottom=2.3mm,
    toptitle=1.25mm, bottomtitle=1.25mm,
    fontupper=\normalsize\linespread{1.26}\selectfont,
    fonttitle=\bfseries\large,
    title={#2\ #3}, #1
}
\newtcolorbox{defBox}[3][]{
    bigbox, colback=DefBg, colframe=DefColor,
    title={#2\ #3}, #1
}
\newtcolorbox{exBox}[3][]{
    bigbox, colback=ExBg, colframe=ExColor,
    title={#2\ #3}, #1
}

\newenvironment{theorem}[1][]{\begin{thmBox}{定理}{#1}}{\end{thmBox}}
\newenvironment{lemma}[1][]{\begin{thmBox}{引理}{#1}}{\end{thmBox}}
\newenvironment{proposition}[1][]{\begin{thmBox}{命题}{#1}}{\end{thmBox}}
\newenvironment{corollary}[1][]{\begin{thmBox}{推论}{#1}}{\end{thmBox}}
\newenvironment{definition}[1][]{\begin{defBox}{定义}{#1}}{\end{defBox}}
\newenvironment{example}[1][]{\begin{exBox}{例题}{#1}\setlength{\parskip}{0.4em}}{\end{exBox}}

% ---------- 盒子：补证（淡底深字，弱标题栏——不跟定理族抢注意力） ----------
\definecolor{SuppColor}{RGB}{150,105,40}
\definecolor{SuppBg}{RGB}{253,250,242}
\definecolor{SuppTitleBg}{RGB}{245,234,210}
\newtcolorbox{suppBox}[3][]{
    bigbox, colback=SuppBg, colframe=SuppColor,
    coltitle=SuppColor, colbacktitle=SuppTitleBg,
    title={#2\ #3}, #1
}
\newenvironment{supplement}[1][]{\begin{suppBox}{补证}{#1}}{\end{suppBox}}

% ---------- 盒子：注（石板蓝，与定理/定义/例区分，弱标题栏） ----------
% 用于原书 Remark 一类「依附于某个定义、但不是定理」的对象
\definecolor{RemColor}{RGB}{74,90,114}
\definecolor{RemBg}{RGB}{247,248,250}
\definecolor{RemTitleBg}{RGB}{231,235,241}
\newtcolorbox{remBox}[3][]{
    bigbox, colback=RemBg, colframe=RemColor,
    coltitle=RemColor, colbacktitle=RemTitleBg,
    title={#2\ #3}, #1
}
\newenvironment{remark}[1][]{\begin{remBox}{注}{#1}}{\end{remBox}}

% ---------- 盒子：夹注（解释 / 结论） ----------
\newtcolorbox{expBox}{
    noteinline, colback=ExpBg, colframe=ExpBg,
    underlay={\draw[ExpColor,line width=2.2pt]
        (frame.south west) -- (frame.north west);}
}
\newtcolorbox{conclBox}{
    noteinline, colback=ConclBg, colframe=ConclBg,
    underlay={\draw[ConclColor,line width=2.2pt]
        (frame.south west) -- (frame.north west);}
}

\newenvironment{explanation}[1][]{%
    \begin{expBox}%
    \setlength{\parskip}{0.45em}\setlength{\parindent}{0pt}%
    \ifstrempty{#1}{}{\textcolor{ExpColor}{\bfseries #1}\par\vspace{1pt}}%
}{%
    \end{expBox}%
}
\newenvironment{conclusion}[1][]{%
    \begin{conclBox}%
    \setlength{\parskip}{0.45em}\setlength{\parindent}{0pt}%
    \ifstrempty{#1}{}{\textcolor{ConclColor}{\bfseries #1}\par\vspace{1pt}}%
}{%
    \end{conclBox}%
}

% ---------- 证明 ----------
% 显式 \small：夹在定理盒（\normalsize 11pt）里时也保持与盒外证明同一字号
\renewenvironment{proof}[1][]{%
    \par\vspace{0.4em}\small\noindent
    \textbf{证明\ifstrempty{#1}{}{（#1）}}\quad\ignorespaces
}{%
    \ifhmode\unskip\nobreak\hfill$\square$\fi
    \ifvmode\noindent\hfill$\square$\fi
    \par\vspace{0.25em}%
}

% ---------- 页眉页脚 ----------
\pagestyle{fancy}
\fancyhf{}
\fancyhead[L]{\small\kaishu\ShortTitle}
\fancyhead[R]{\small\kaishu\DocSubTitle}
\fancyfoot[C]{\small\thepage}
\renewcommand{\headrulewidth}{0.4pt}
\renewcommand{\headrule}{\hbox to\headwidth{\color{black!25}\leaders\hrule height \headrulewidth\hfill}}

\fancypagestyle{plain}{
    \fancyhf{}
    \fancyhead[L]{\small\kaishu\ShortTitle}
    \fancyhead[R]{\small\kaishu\DocSubTitle}
    \fancyfoot[C]{\small\thepage}
    \renewcommand{\headrulewidth}{0.4pt}
    \renewcommand{\headrule}{\hbox to\headwidth{\color{black!25}\leaders\hrule height \headrulewidth\hfill}}
}

% ---------- 封面 ----------
\newcommand{\MakeTitlePage}{%
    \begin{titlepage}
        \centering
        \vspace*{4.5cm}
        {\fontsize{30pt}{38pt}\selectfont\bfseries\color{ThmColor}\DocMainTitle\par}
        \vspace{1.1cm}
        {\fontsize{16pt}{22pt}\selectfont\DocSubTitle\par}
        \vfill
        {\large\kaishu\DocSource\par}
        \vspace{0.7cm}
        {\large\DocDate\par}
        \vspace{2.0cm}
    \end{titlepage}%
}

% ---------- 链接与书签 ----------
\usepackage[unicode,colorlinks=true,linkcolor=ThmColor,citecolor=ExColor,
            urlcolor=ConclColor,bookmarksnumbered=true,bookmarksopen=true]{hyperref}
\hypersetup{
    pdftitle={\DocMainTitle\ \DocSubTitle},
    pdfauthor={\DocSource},
    pdfcreator={XeLaTeX},
}

\numberwithin{equation}{chapter}
\setcounter{secnumdepth}{-1}
\setcounter{tocdepth}{2}
\xeCJKDeclareCharClass{CJK}{"2460 -> "2473}

% ---------- 全局字号：正文档 10pt，定理族 11pt（2026-11-03 三次调整 9pt→10pt） ----------
% 定理族盒内显式 \normalsize（11pt）；其余全部 \small（10pt）——
% 正文、定义、例题、证明、解释、结论、页眉页脚用同一个字号，只靠颜色与盒形区分层次。
% 注意 \small 是绝对映射（11pt 类中恒为 10pt），重复套用不会继续缩小；
% 改档时 proof 与 noteinline 里的显式字号必须同步（共三处：AtBeginDocument、proof、noteinline）。
\AtBeginDocument{\small}

%==========================================================
%  以下由 md2tex.py 拼接
%==========================================================
